\section{Introduction}
\label{sec:introduction}

\subsection{The leave-one-out stability trap}

Practitioners often read low leave-one-out cross-validation error as evidence
that a local classifier is not overly sensitive to any single training point.
The reasoning seems natural: if removing a point barely changes the prediction
at that same point, the classifier must be stable.  This paper argues that the
inference is invalid for deterministic $k$-nearest-neighbor classification, and
that the flaw is conceptual, not asymptotic.

The central question is simple but easy to overlook:

\begin{quote}
When a perturbation is applied to the training sample, what exactly has been
perturbed, and where is the perturbed predictor evaluated?
\end{quote}

Leave-one-out (LOO) couples these two choices: the perturbation is delete-one,
and the evaluation point is fixed to the deleted occurrence.  Replacement
robustness (replace-one), in contrast, decouples them: any sample occurrence
may be replaced by any labeled point, and the evaluation query is arbitrary.
This paper shows that this difference in \emph{perturbation semantics} produces
a sharp structural separation between the two notions for deterministic $k$-NN,
for every $k \ge 1$.

\subsection{A two-point preview}

Before developing the general machinery, we give the simplest possible witness.
Consider a metric space with two vertices $a,b$ joined by a single edge, and
the ordered sample
\[
S = \big((a,0),\;(b,0)\big).
\]

\paragraph{LOO is zero.}
Delete the first occurrence $(a,0)$ and evaluate at $a$.  The sole remaining
neighbor $(b,0)$ votes $0$, so the loss at $(a,0)$ is unchanged.  The same
holds for the second occurrence.  Hence $\Delta_{\mathrm{loo}}(S,i)=0$ for
$i=1,2$.

\paragraph{Replace-one flips.}
Now replace the first occurrence with $(a,1)$.  The perturbed sample is
$((a,1),(b,0))$.  At query $a$, the nearest neighbor is now $(a,1)$, so the
prediction flips from $0$ to $1$.  The loss at $(a,0)$ changes from $0$ to
$1$, giving $\Delta_{\mathrm{rep}}=1$.

The separation already appears in a two-point metric space with the simplest
possible classifier.  No asymptotics, randomness, or high-dimensional geometry
is needed.  The rest of the paper generalises this mechanism to all $k \ge 1$.

\subsection{Contributions}

\begin{enumerate}[leftmargin=*,label=(\arabic*)]
    \item \textbf{Perturbation semantics.}  We formalize delete-one,
        insert-one, replace-one, and leave-one-out as distinct operations
        coupled to distinct evaluation-point conventions
        (\Cref{sec:perturbation-semantics}).  The key distinction is that LOO
        is not simply delete-one: it is delete-one with the evaluation point
        constrained to be the deleted occurrence.
    \item \textbf{Uniform perturbation calculus.}  Replace-one decomposes into
        delete-one followed by insert-one, yielding the bound
        $\beta_{\mathrm{rep}}(n) \le \beta_{\mathrm{del}}(n) +
        \beta_{\mathrm{ins}}(n-1)$ at the uniform level
        (\Cref{sec:uniform-calculus}).  LOO cannot substitute for delete-one
        control in this bound.
    \item \textbf{Margin mechanism and separations.}  For deterministic $k$-NN,
        a prediction changes if and only if the signed top-$k$ vote margin
        crosses zero (\Cref{sec:margin-mechanism}).  Using this criterion, we
        construct LOO-stable but replace-one-unstable samples for every odd $k
        \ge 1$ (\Cref{sec:odd-k-separation}) and for even $k$ under
        deterministic tie-breaking (\Cref{sec:even-k-separation}).
    \item \textbf{Diagnostic and empirical evidence.}  A margin-based
        diagnostic ($|\margin_k(S,x)| \le 2$) detects queries vulnerable to
        single-replacement flips.  Experiments on synthetic graph metrics and
        UCI tabular data show that the LOO/replace-one gap is not confined to
        hand-crafted witnesses but appears at measurable frequency
        (\Cref{sec:experiments}).
\end{enumerate}

\noindent The paper is not a critique of leave-one-out cross-validation as a
practical tool.  It is a precise statement about the semantic boundary of what
LOO can certify: LOO is a coupled deletion diagnostic, not a certificate of
replacement robustness.
